\documentclass[11pt,a4paper]{article}
\usepackage[T1]{fontenc}
\usepackage{lmodern,amsmath,amssymb,amsthm,mathtools,microtype}
\usepackage[margin=27mm]{geometry}
\usepackage[colorlinks=true,linkcolor=blue,citecolor=blue,urlcolor=blue]{hyperref}
\newtheorem{theorem}{Theorem}[section]
\newtheorem{lemma}[theorem]{Lemma}
\newtheorem{corollary}[theorem]{Corollary}
\newcommand{\Q}{\mathbb Q}
\DeclareMathOperator{\Nm}{N}
\setlength{\emergencystretch}{2em}
\title{Three quadratic equations for two translated elliptic parameters}
\author{}\date{10 November 2026}
\begin{document}\maketitle
\begin{abstract}
Let $E:y^2=x^3-Dx$ over a field of characteristic zero and let
$P=(a,b)$ be a fixed point satisfying the specified nonvanishing
conditions. We give three quadratic equations in four field coordinates
whose solutions describe exactly the points $H$ for which $H$ and
$H-P$ are finite with nonzero abscissas. The coordinates retain both
reciprocal abscissas and both values of $y/x$. All exceptional leading
coefficients are treated explicitly. The resulting direct equations
for formal-parameter ratios save one common field coordinate and one
field coordinate for each numerator pair. Every identity has a separate
exact verification using only the Python standard library.
\end{abstract}

\section{The quadratic equations}
Let $k$ be a field of characteristic zero, let $D\in k^\times$, and put
\[
 E:y^2=x^3-Dx,\qquad P=(a,b)\in E(k).
\]
Assume
\begin{equation}\label{eq:assumptions}
 ab\ne0,\qquad C:=D^2+6a^2D-3a^4\ne0.
\end{equation}
Write $A=a^2+D$ and $B=a^2-D$. The point equation is $b^2=aB$.
For four field coordinates $S,T,f,g$, define
\begin{align*}
 L_f&=2bD(aT-b),& Q_f&=aAST+bAS-bBT+4a^2D,\\
 L_g&=2bD(aS+b),& Q_g&=-aAST-bBS+bAT-4a^2D.
\end{align*}
Consider the three quadratic equations
\begin{equation}\label{eq:system}
 L_ff=Q_f,\qquad L_gg=Q_g,\qquad
 D^2fg+aD(f+g)+aST+b(S-T)+D=0.
\end{equation}
No inverse witness occurs in this system.

\begin{theorem}\label{thm:chart}
Equations~\eqref{eq:system} describe exactly the points
\[
 H\in E(k)\setminus\{O,(0,0),P,P+(0,0)\}
\]
through the identities
\begin{equation}\label{eq:interpretation}
 S=\frac{y(H)}{x(H)},\quad T=\frac{y(H-P)}{x(H-P)},\quad
 f=\frac1{x(H)},\quad g=\frac1{x(H-P)}.
\end{equation}
In particular every solution has $fg\ne0$. Every allowed point has
one and only one quadruple of coordinates satisfying these identities,
and every solution of the quadratic system is such a quadruple.
The theorem also holds when $a^2+D=0$.
\end{theorem}

\section{Polynomial reconstruction and its exceptional cases}
Define
\[
 \Phi=aS^2T^2-B(S^2+T^2)+2AST+4aD
\]
and two polynomial expressions in $S,T$:
\begin{align}
 2bDf_0&=-AS+BT-bS^2-aS^2T,\label{eq:reconstruct}\\
 2bDg_0&=AT-BS-bT^2+aST^2.\nonumber
\end{align}
The divisions here are only by the fixed nonzero constant $2bD$.
The following identities hold as polynomial identities modulo
$b^2-a^3+aD$:
\begin{align}
 Df_0^2+S^2f_0-1
 &=\frac{a^2S^2-b^2}{4ab^2D}\Phi,\label{eq:curve}\\
 Dg_0^2+T^2g_0-1
 &=\frac{a^2T^2-b^2}{4ab^2D}\Phi,\nonumber\\
 L_ff_0-Q_f&=-a\Phi,\qquad L_gg_0-Q_g=a\Phi.\label{eq:quotients}
\end{align}
Set
\[
 r_f=L_ff-Q_f,\qquad r_g=L_gg-Q_g,
 \qquad R=D^2fg+aD(f+g)+aST+b(S-T)+D.
\]
The elimination identity is
\begin{equation}\label{eq:master}
\begin{split}
 L_fL_gR={}&D^2r_fr_g+(aDL_g+D^2Q_g)r_f\\
          &+(aDL_f+D^2Q_f)r_g-aD^2C\Phi.
\end{split}
\end{equation}
One also has
\begin{equation}\label{eq:compatibility}
 (2bD)^2R(f_0,g_0)=D^2\bigl(B-aST-b(S-T)\bigr)\Phi.
\end{equation}
All these identities follow by expansion and substitution of
$b^2=a(a^2-D)$. In particular, no division by $L_f$ or $L_g$ is
used in~\eqref{eq:master}.

\begin{lemma}\label{lem:unique}
Every solution of~\eqref{eq:system} satisfies $\Phi=0$ and
$f=f_0$, $g=g_0$.
\end{lemma}
\begin{proof}
The nonvanishing of $aD^2C$ and~\eqref{eq:master} give $\Phi=0$.
When $L_fL_g\ne0$, equations~\eqref{eq:quotients} give the claim.
Suppose $L_f=0$. Then $T=b/a$. If $A=0$, substitution gives
$\Phi=-8a^3\ne0$, a contradiction. If $A\ne0$, the equation
$\Phi=0$ is linear in $S$ and gives
\[
 S=\frac{a^4-6Da^2+D^2}{2bA},\qquad
 L_g=-\frac{aDC}{A}\ne0.
\]
Thus $g=g_0=1/a$. The coefficient of $f$ in $R$ is
$D^2/a+aD=DA/a\ne0$, and~\eqref{eq:compatibility} then gives
$f=f_0$. Interchanging $S$ and $T$ and replacing $b$ by $-b$
handles $L_g=0$. The two leading coefficients cannot vanish
together: at $S=-b/a,T=b/a$, one has $Q_f=C\ne0$.
\end{proof}

\section{Identification with elliptic translation}
By~\eqref{eq:curve}, the equation $\Phi=0$ implies $f_0g_0\ne0$.
It therefore supplies points
\[
 H=(f_0^{-1},S/f_0),\qquad H'=(g_0^{-1},T/g_0)
\]
on $E$. Two further identities are
\begin{equation}\label{eq:line}
 a(f_0T+g_0S)+b(g_0-f_0)=S+T
\end{equation}
and
\begin{equation}\label{eq:abscissa}
\begin{split}
 f_0g_0(aS+b)^2-a(1-af_0)^2
 =-\frac{(aS+b)^2\bigl(aST+b(S-T)+D\bigr)}{4b^2D^2}\Phi.
\end{split}
\end{equation}
The first is an identity even before imposing $\Phi=0$.

If $x=x(H)\ne a$, the slope of the chord through $H$ and $-P$ is
\[
 \mu=\frac{xS+b}{x-a}.
\]
The curve equation $x^2-S^2x-D=0$ and $b^2=a(a^2-D)$ give
\[
 \mu^2-x-a=\frac{x(aS+b)^2}{a(x-a)^2}.
\]
Equation~\eqref{eq:abscissa} says that this is $g_0^{-1}=x(H')$.
After multiplying by $f_0g_0$, equation~\eqref{eq:line} is precisely
the collinearity condition for $H,-H',-P$. It therefore identifies the
ordinate too, and proves $H'=H-P$.

Suppose $x(H)=a$. Then $S=\pm b/a$. At $S=b/a$ the expression
for $f_0$ is $-a/D$. If $A\ne0$, this differs from $1/a$, so the
case is impossible. If $A=0$, substitution in $\Phi$ gives
$-8a^3\ne0$, which is again impossible. At $S=-b/a$, the same
$A=0$ calculation excludes the case; otherwise
\[
 f_0=1/a,\quad
 T=-\frac{a^4-6Da^2+D^2}{2bA},\quad
 g_0=\frac{4b^2}{A^2}.
\]
The doubling formulas identify $H=-P$ and $H'=-2P$.
This proves the translation assertion for every solution.

Conversely, let $H$ and $H-P$ be finite with nonzero abscissas.
For $x(H)\ne a$, substitution of the chord formulas gives $\Phi=0$
and~\eqref{eq:reconstruct}. The preceding explicit doubling calculation
covers $H=-P$; the other possibility $H=P$ is outside the allowed
set. Equations~\eqref{eq:quotients} and~\eqref{eq:compatibility}
then imply~\eqref{eq:system}. Uniqueness follows from the four
rational functions in~\eqref{eq:interpretation}. This proves
Theorem~\ref{thm:chart}.

\section{Necessity of an exceptional-constant condition}
The condition $C\ne0$ cannot simply be dropped. Over
$k=\Q(\sqrt3)$, take
\[
 a=1,\qquad D=-3+2\sqrt3,\qquad b=\sqrt3-1.
\]
Then $abD\ne0$, $b^2=1-D$, and $C=0$. Set
\[
 S=-b,\qquad T=b,\qquad f=0,\qquad g=3/D-4.
\]
Here $L_f=L_g=Q_f=Q_g=0$, while the third equation reduces to
$Dg+4D-3=0$. Thus this is a solution of the three quadrics with
$f=0$, and it cannot have the interpretation~\eqref{eq:interpretation}.
The elimination identity identifies the precise constant responsible
for this failure.

\section{Direct equations for formal-parameter ratios}
Write $z=-x/y$. Suppose numerator points have a quadratic description
with coordinates $V,Q,I,K$ such that
\[
 z(U)=-V/I,\qquad z(V_1)=-Q/K,\qquad IK\ne0.
\]
Zero numerators may represent the identity. At a common-point solution
with $ST\ne0$, the two ratios are given directly by
\[
 LI=VS,\qquad L'K=QT.
\]
These are quadratic equations. They have exactly the values
$L=z(U)/z(H)$ and $L'=z(V_1)/z(H-P)$, including either identity
numerator. Consequently no separate coordinate for $z(V_1)$ is needed.

Over a number field of degree $d$, expansion in a fixed rational basis
turns each field coordinate into $d$ rational coordinates and each
quadratic field equation into $d$ rational quadratic equations.
The common-point model uses $4d$ coordinates and $3d$ equations.
In a construction with $r$ numerator pairs, replacing a five-coordinate
common chart and eliminating one numerator parameter per pair saves
$d(r+1)$ rational coordinates, provided the construction already
supplies the stated domains and the numerator identities.
This is an algebraic statement about representations; it does not
supply any Diophantine definition of a ring of integers by itself.

For the applications over $\Q(\rho)$, $\rho^2=2$, the point is
\[
 a=\frac{57-28\rho}{8},\qquad b=\frac{-660+417\rho}{32},
 \qquad D=11-6\rho.
\]
Exact calculation gives
\[
 C=-\frac{94897379}{4096}+\frac{4195821}{256}\rho,
 \qquad \Nm(C)=-\frac{8203357319351}{16777216}\ne0.
\]
Thus the common-chart saving is two rational coordinates, and each
numerator pair saves two more. All other arithmetic arguments remain
those of the corresponding defining-polynomial construction.

\section{Exact verification}
The program
\begin{center}\small
\nolinkurl{computation/elliptic_constraints/check_fixed_difference_chart.py}
\end{center}
verifies twelve general identities in the quotient ring
\[
 \mathbb Z[a,b,D,S,T]/(b^2-a^3+aD).
\]
It uses explicit integer polynomial arithmetic, imports no producer,
and checks the master identity without division. It also checks the
curve identities, translation, reconstruction, and the nonzero
coefficients in the exceptional cases. Independent rational-point
checks compare the formulas with the chord-and-tangent law at
$H=-P,2P,3P$ on $y^2=x^3-2x$ with $P=(2,2)$.
Its output is
\nolinkurl{fixed-difference-chart-certificate.json}.
The separate quartic checker verifies the displayed nonzero constant
in $\Q(\sqrt2)$ and the resulting rational component equations.
The identities prove the assertions for every allowed field and point;
the sample points serve only as additional verification.
\end{document}
